% Generated from locale/tr/content/first-order-logic/syntax-and-semantics/intro-syntax.tex; do not hand-edit.


\olfileid[tr]{fol}{syn}{itx}

\olsection{Giriş}

Birinci dereceden mantığın kuramını ve metakuramını geliştirebilmek için
önce ifadelerinin sözdizimini ve semantiğini tanımlamalıyız. Birinci
dereceden mantığın ifadeleri terimlerden ve !!{formula}s{}den oluşur.
Terimler !!{variable}s{}den, !!{constant}s{}den ve !!{function}s{}dan
oluşturulur. !!^{formula}s ise !!{predicate}s ile terimlerden (bunlar en
küçük, ``atomik'' !!{formula}s{}dir), ardından da mantıksal
bağlaçlar ve niceleyiciler kullanılarak atomik !!{formula}s{}den daha
karmaşık formüller oluşturmak suretiyle kurulur. Oluşum kurallarını
ortaya koymanın birçok farklı yolu vardır; burada bunlardan yalnızca
birini veriyoruz. Başka sistemler farklı simgeler seçebilir, farklı
bağlaç kümelerini ilkel kabul edebilir ve parantezleri farklı biçimde
kullanabilir (hatta Polonya gösterimi denen yöntemde olduğu gibi hiç
kullanmayabilir). Yine de bütün yaklaşımların ortak yanı, terimler ve
!!{formula}s kümesini oluşum kurallarıyla \emph{tümevarımsal olarak}
tanımlamalarıdır. Tanım doğru kurulmuşsa her ifade, oluşum kurallarına
göre özünde yalnızca bir yoldan elde edilebilir. İfadelerin \emph{tek
türlü okunabilir} olmasını sağlayan bu tümevarımsal tanım, aynı
yöntemle---tümevarımsal tanımla---bu ifadelere anlam vermemize olanak
tanır.

